\documentclass[11pt]{article}
\usepackage[a4paper,margin=2.2cm]{geometry}
\usepackage{amsmath,amssymb}
\usepackage[T1]{fontenc}
\usepackage{lmodern}
\usepackage{xcolor}
\usepackage{listings}
\usepackage{titlesec}
\usepackage{parskip}
\definecolor{codebg}{gray}{0.95}
\definecolor{notegray}{gray}{0.40}
\lstset{basicstyle=\ttfamily\small, backgroundcolor=\color{codebg}, frame=single, rulecolor=\color{gray!40},
        breaklines=true, breakatwhitespace=true, columns=fullflexible, keepspaces=true, upquote=true, aboveskip=2pt, belowskip=8pt}
\titleformat{\section}{\large\bfseries}{}{0pt}{}
\titlespacing*{\section}{0pt}{14pt}{4pt}
\newcommand{\codenote}[1]{{\small\color{notegray}\emph{Code:} #1\par}}
\begin{document}
{\LARGE\bfseries Statistics in the turbulence talk figures}\par\smallskip
{\small Formulas recovered from \texttt{codes/talk\_figures.py}, \texttt{codes/paper\_figures\_li.py},\\
\texttt{turbulence/paper\_figs\_ensemble.py}, \texttt{codes/check\_moments.py}.\\
Each formula is followed by its LaTeX code (copy from the grey box; math mode, no custom macros).
Grey \emph{Code:} lines are implementation details, not for the slides.}

\section{Notation (common to all figures)}
Signals $x$ of $T$ samples $x(t)$, $t = 0,\dots,T-1$ ($N$ signals). Lags $\tau$ are integers (samples). Lagrangian figures show the lag as $\tau/\tau_\eta$ ($\tau_\eta$: Kolmogorov time in samples); Eulerian (jets) figures show $\tau$ in samples.

\[
\delta_\tau x(t) = x(t+\tau) - x(t)
\]
\begin{lstlisting}
\delta_\tau x(t) = x(t+\tau) - x(t)
\end{lstlisting}
$\mathbb{E}_x$: average over signals. $\mathbb{E}_t$: average over time. $\mathbb{E}_{x,t}$: both. $\mathrm{Var}$, $\mathrm{Std}$, $\mathrm{Median}$, $\mathrm{Quantile}$ carry the same subscripts. All are empirical (over the signals shown).

\section{Structure functions $S_p(\tau)$}
Figures: \texttt{li\_structure\_functions} ($p = 2,4,6,8$, Lagrangian); \texttt{fig6\_*\_1}, right panel ($p = 4,6,8$, both). Curves shifted down by $10^{-1}$ per order.

\[
S_p(\tau) = \mathbb{E}_{x,t}\big[ (\delta_\tau x)^p \big]
\]
\begin{lstlisting}
S_p(\tau) = \mathbb{E}_{x,t}\big[ (\delta_\tau x)^p \big]
\end{lstlisting}
\codenote{Fig.~6 takes the absolute value of the mean, $|\mathbb{E}_{x,t}[(\delta_\tau x)^p]|$: identical for even $p$. Lags: $\tau = 1,\dots,T/2-1$ (Fig.~6), $\tau = 1,\dots,\min(\texttt{max\_lag}, T-2)$ (Li).}
\section{Flatness $F_p(\tau)$}
Figure: \texttt{li\_flatness} ($p = 4,6,8$, Lagrangian).

\[
F_p(\tau) = \frac{S_p(\tau)}{S_2(\tau)^{p/2}}
\]
\begin{lstlisting}
F_p(\tau) = \frac{S_p(\tau)}{S_2(\tau)^{p/2}}
\end{lstlisting}
Gaussian values, for reference (not drawn in the figure):

\[
F_4 = 3, \qquad F_6 = 15, \qquad F_8 = 105
\]
\begin{lstlisting}
F_4 = 3, \qquad F_6 = 15, \qquad F_8 = 105
\end{lstlisting}
\section{Flatness slope}
Figure: \texttt{fig6\_*\_2}, left panel (both). Dotted line at $0$.

\[
\frac{d \log F_4}{d \log \tau}
\]
\begin{lstlisting}
\frac{d \log F_4}{d \log \tau}
\end{lstlisting}
\codenote{Evaluated on 30 log-spaced lags in $[1, T/2)$, centred second-order finite differences in $\log\tau$ (\texttt{np.gradient}).}
\section{Local scaling exponent $\zeta_p(\tau)$}
Figure: \texttt{li\_local\_exponents} ($p = 4,6,8$, Lagrangian). Dotted line at $p/2$ (no intermittency: $S_p \propto S_2^{\,p/2}$).

\[
\zeta_p(\tau) = \frac{d \log S_p}{d \log S_2}
\]
\begin{lstlisting}
\zeta_p(\tau) = \frac{d \log S_p}{d \log S_2}
\end{lstlisting}
\codenote{Computed as the ratio $\frac{d \log S_p / d \log \tau}{d \log S_2 / d \log \tau}$ (same thing, chain rule), finite differences on the unit-spaced lag grid $\tau = 1, 2, \dots$, as Li et al.\ 2024.}
\section{Increment PDFs}
\[
\mathrm{PDF}\!\left( \frac{\delta_\tau x}{\mathrm{Std}_{x,t}[\delta_\tau x]} \right)
\]
\begin{lstlisting}
\mathrm{PDF}\!\left( \frac{\delta_\tau x}{\mathrm{Std}_{x,t}[\delta_\tau x]} \right)
\end{lstlisting}
\texttt{li\_increment\_pdfs} (Lagrangian): normalised density, $\tau/\tau_\eta = 1, 2, 5$, shifted down by $10^{-3}$ per lag.\\ \texttt{fig6\_*\_2}, middle panel (both): PDF divided by its maximum (peak $=1$), $\tau = 1, 4, 16, 64$ samples, shifted down by $10^{-1}$ per lag.

\codenote{Fig.~6: histogram on $[-6, 6]$, 201 bin edges, Gaussian smoothing of 1.5 bins. Li: no smoothing, bins over $\pm\max|\delta_\tau x|/\mathrm{Std}_{x,t}[\delta_\tau x]$.}
\section{Marginal PDF}
Figure: \texttt{fig6\_*\_1}, left panel (both).

\[
\mathrm{PDF}\!\left( \frac{x}{\mathrm{Std}_{x,t}[x]} \right)
\]
\begin{lstlisting}
\mathrm{PDF}\!\left( \frac{x}{\mathrm{Std}_{x,t}[x]} \right)
\end{lstlisting}
\codenote{The data are standardised ($\mathbb{E}_{x,t}[x] = 0$, $\mathrm{Std}_{x,t}[x] = 1$), so the code histograms $x$ directly; the model samples are put on the same bins without re-standardising. Gaussian smoothing of 1.5 bins.}
\section{Power spectrum $P(\omega)$}
Figure: \texttt{fig6\_*\_1}, middle panel (both); $\omega = k/T$ ($k = 1,\dots,T/2$) in cycles per sample, shown for $0 < \omega \le 0.4$.

\[
P(\omega) = \mathbb{E}_x \Big[ \Big| \sum_{t=0}^{T-1} x(t)\, e^{-2\pi i \omega t} \Big|^2 \Big]
\]
\begin{lstlisting}
P(\omega) = \mathbb{E}_x \Big[ \Big| \sum_{t=0}^{T-1} x(t)\, e^{-2\pi i \omega t} \Big|^2 \Big]
\end{lstlisting}
\codenote{No $1/T$ normalisation (\texttt{|rfft|}$^2$ averaged over signals).}
\section{Skewness}
Figure: \texttt{fig6\_*\_2}, right panel (both). Dotted line at $0$.

\[
\mathrm{skewness}(y) = \mathbb{E}_{x,t}\left[ \left( \frac{y - \mathbb{E}_{x,t}[y]}{\mathrm{Std}_{x,t}[y]} \right)^{3} \right], \qquad y = x(t)\,\delta_\tau x(t)
\]
\begin{lstlisting}
\mathrm{skewness}(y) = \mathbb{E}_{x,t}\left[ \left( \frac{y - \mathbb{E}_{x,t}[y]}{\mathrm{Std}_{x,t}[y]} \right)^{3} \right], \qquad y = x(t)\,\delta_\tau x(t)
\end{lstlisting}
\codenote{\texttt{skew\_mode=\textquotesingle power\textquotesingle} (the talk setting). 45 log-spaced lags in $[1, T/2)$.}
\section{Cross-scale correlations $C_{p,q}(\tau, \tau^\star)$}
Figure: \texttt{fig6\_*\_3} (both), $\tau^\star = 20$ samples (dotted vertical line; shown as $\tau^\star/\tau_\eta$ in the Lagrangian figure). Horizontal line at $1$: value when the increment amplitudes at the two scales are uncorrelated.

\[
C_{p,q}(\tau, \tau^\star) = \frac{\mathbb{E}_{x,t}\big[ |\delta_{\tau^\star} x|^p \, |\delta_\tau x|^q \big]}{\mathbb{E}_{x,t}\big[ |\delta_{\tau^\star} x|^p \big] \, \mathbb{E}_{x,t}\big[ |\delta_\tau x|^q \big]}
\]
\begin{lstlisting}
C_{p,q}(\tau, \tau^\star) = \frac{\mathbb{E}_{x,t}\big[ |\delta_{\tau^\star} x|^p \, |\delta_\tau x|^q \big]}{\mathbb{E}_{x,t}\big[ |\delta_{\tau^\star} x|^p \big] \, \mathbb{E}_{x,t}\big[ |\delta_\tau x|^q \big]}
\end{lstlisting}
The three panels are $C_4 \equiv C_{2,2}$, \; $C_6 \equiv C_{3,3}$ and $C_{4,2}$. For even $p, q$ the absolute values are not needed, e.g.

\[
C_4(\tau, \tau^\star) = \frac{\mathbb{E}_{x,t}\big[ (\delta_{\tau^\star} x)^2 \, (\delta_\tau x)^2 \big]}{S_2(\tau^\star)\, S_2(\tau)}
\]
\begin{lstlisting}
C_4(\tau, \tau^\star) = \frac{\mathbb{E}_{x,t}\big[ (\delta_{\tau^\star} x)^2 \, (\delta_\tau x)^2 \big]}{S_2(\tau^\star)\, S_2(\tau)}
\end{lstlisting}
\codenote{Both increments start at the same time $t$; the numerator averages over $t \le T-1-\max(\tau, \tau^\star)$, each denominator factor over its own full range. $10^{-8}$ added to the denominator. Lags $\tau = 1,\dots,T/2-1$.}
\section{Rare signals: acceleration criterion}
Acceleration = one-step increment (as Li et al.):

\[
a(t) = \delta_1 x(t) = x(t+1) - x(t)
\]
\begin{lstlisting}
a(t) = \delta_1 x(t) = x(t+1) - x(t)
\end{lstlisting}
Each signal is summarised by its peak acceleration $\max_t |a(t)|$. \textbf{Criterion of the talk}: a signal is rare if

\[
\max_t |a(t)| \;>\; \mathrm{Quantile}_{0.95,\;x\in\mathrm{data}}\big[ \max_t |a(t)| \big]
\]
\begin{lstlisting}
\max_t |a(t)| \;>\; \mathrm{Quantile}_{0.95,\;x\in\mathrm{data}}\big[ \max_t |a(t)| \big]
\end{lstlisting}
i.e.\ rare = the $5\,\%$ of the data signals with the largest peak acceleration.

\textbf{Alternative} (\texttt{turb\_rare\_typical\_sigma.ipynb}, threshold in units of the standard deviation of $a$): rare if

\[
\max_t |a(t)| \;\ge\; 20\,\mathrm{Std}_{x,t\,;\,x\in\mathrm{data}}[a]
\]
\begin{lstlisting}
\max_t |a(t)| \;\ge\; 20\,\mathrm{Std}_{x,t\,;\,x\in\mathrm{data}}[a]
\end{lstlisting}
\codenote{One single number (every time of every data signal), so the threshold is the same level for the data and for the generated signals.}
\codenote{\texttt{rare\_events\_data\_vs\_generated} applies the same rule at the lags $\tau/\tau_\eta = 20, 5, 1$, with $a$ replaced by $\delta_\tau x$ and the quantile taken at that lag on the data; the generated signals use the same absolute threshold. \texttt{rare\_events\_by\_logp} instead picks signals among the $1\,\%$ with the highest $-\log p$.}
\section{$-\log p$: what the mean and the std are taken over}
The model is $p(x) \propto \exp\!\big(+\Theta \cdot \phi(x)\big)$ ($\Theta = \Theta_{\rm reg}$ at the last time step), so

\[
-\log p(x) = -\,\Theta \cdot \phi(x) + \mathrm{const}
\]
\begin{lstlisting}
-\log p(x) = -\,\Theta \cdot \phi(x) + \mathrm{const}
\end{lstlisting}
The constant is unknown, so the level of $-\log p$ means nothing: only differences between signals do. Large $-\log p$ = unlikely signal (rare signals sit on the right of the histograms). One fixed $\Theta$; the $N = 10000$ data signals are the training signals:

\[
\mathrm{mean} = \mathbb{E}_{x\in\mathrm{data}}\big[ -\log p(x) \big], \qquad \mathrm{std} = \mathrm{Std}_{x\in\mathrm{data}}\big[ -\log p(x) \big]
\]
\begin{lstlisting}
\mathrm{mean} = \mathbb{E}_{x\in\mathrm{data}}\big[ -\log p(x) \big], \qquad \mathrm{std} = \mathrm{Std}_{x\in\mathrm{data}}\big[ -\log p(x) \big]
\end{lstlisting}
\begin{itemize}
\item \textbf{Table} (per category $c$ of potentials): the same $\mathbb{E}_{x\in\mathrm{data}}$ and $\mathrm{Std}_{x\in\mathrm{data}}$, applied to the contribution $-\sum_{F \in c} \Theta_F \cdot \phi_F(x)$ instead of $-\log p(x)$. Computed for each seed ($\Theta$ differs from seed to seed, the data signals are the same), then averaged over the seeds: $\mathbb{E}_{\rm seed}[\,\cdot\,]$.
\item \textbf{Bar plots}: same quantities; error bars = $\mathrm{Std}_{\rm seed}$. Third panel: $\mathrm{Std}_{x}[\text{category}] \,/\, \mathrm{Std}_{x}[-\log p]$.
\item \textbf{Histograms} (data signals, or generated signals, each with its own median): $-\log p(x) - \mathrm{Median}_{x}\big[-\log p(x)\big]$.
\item \textbf{Gaussian test}: $-\log p(x) - \mathbb{E}_{x\in\mathrm{train}}\big[-\log p(x)\big]$ for the exact $-\log p$, $-\log p(x) - \mathbb{E}_{x\in\mathrm{gen}}\big[-\log p(x)\big]$ for the MGD one. For $M$ dimensions the exact $-\log p = \chi^2_M/2 + \mathrm{const}$, hence $\mathrm{Std} = \sqrt{M/2}$.
\end{itemize}

\codenote{$\mathbb{E}_{x\in\mathrm{data}}[-\log p]$ is a cross-entropy of the model on the data, not the entropy of the model (that would be the mean over samples of $p$). Code: \texttt{ddof=1} for every std ($N/(N-1) = 1.0001$ at $N = 10000$, so $\mathrm{Std}$ above is the same number to four digits). The sign flip leaves every std unchanged and changes the sign of every mean. Category contributions add up to $-\log p + \mathrm{const}$, but their stds do not add up (covariances between categories). The mean of a category contribution is a level only: a constant added to $\phi_F$ shifts it without changing $p$. \texttt{energy\_table}, \texttt{plot\_energy\_bars}, \texttt{plot\_logp\_histogram}, \texttt{plot\_logp\_gaussian} in \texttt{codes/talk\_figures.py}.}
\section{Error bands (all figures)}
Line: the statistic on all signals. Band, \texttt{band=\textquotesingle batch\textquotesingle} (default): min/max of the same statistic over 10 disjoint batches of signals (data and model, as Li et al.). \texttt{band=\textquotesingle seeds\textquotesingle}: model line = mean over independent runs, band = $\pm 1$ std; the data band stays min/max over batches.

\end{document}
